% Generated draft: requires mathematical and source review.
\documentclass[11pt]{article}
\usepackage[margin=25mm]{geometry}
\usepackage{fontspec}
\setmainfont{Latin Modern Roman}
\usepackage{amsmath,amssymb,mathtools,unicode-math}
\setmathfont{Latin Modern Math}
\usepackage{enumitem,needspace,xurl,hyperref}
\setlength{\parindent}{0pt}
\setlength{\parskip}{5pt}
\setlength{\emergencystretch}{3em}
\newcommand{\N}{\mathbb{N}}
\newcommand{\Z}{\mathbb{Z}}
\newcommand{\Q}{\mathbb{Q}}
\newcommand{\R}{\mathbb{R}}
\newcommand{\C}{\mathbb{C}}
\newcommand{\F}{\mathbb{F}}
\newcommand{\A}{\mathbb{A}}
\newcommand{\PP}{\mathbb{P}}
\newcommand{\E}{\mathbb{E}}
\newcommand{\eps}{\varepsilon}
\newcommand{\dd}{\,\mathrm{d}}
\begin{document}
\section*{UnsolvedMath ID: 20000401}
\textbf{Source problem code:} catalogue entry
\subsection{Definitions and mathematical statement}
Jordan Ellenberg's AIM question asks for a geometric interpretation of a variance expression printed as \(\operatorname{variance}\Lambda(n)\sim q^{?}\,n\) as \(q\to\infty\). The supplied report says the exponent question mark is present in the official source, not merely an OCR loss. The averaging ensemble, normalization, and meaning of \(\Lambda(n)\) are unspecified, so the exponent cannot be silently filled in.

The report proposes a common representation-theoretic synthesis for two interpretations. For monic \(f\in\F_q[T]\), define \(\Lambda(f)=\deg P\) if \(f=P^a\) for a monic irreducible \(P\), and zero otherwise. On the squarefree degree-\(n\) configuration space this is \(n\) for irreducible polynomials and zero otherwise. Let \(V_{\mathrm{std}}\) be the \((n-1)\)-dimensional standard representation of \(S_n\) and put
\[
L_n=\lambda_{-1}(V_{\mathrm{std}})
=\sum_{j=0}^{n-1}(-1)^j\bigwedge^jV_{\mathrm{std}}.
\]
The report asserts that its virtual character is \(n\) on \(n\)-cycles and zero on other conjugacy classes, hence realizes \(\Lambda\) on squarefree Frobenius configurations. The summands correspond to the hooks \((n-j,1^j)\).

The claimed raw and centered global leading coefficients are respectively
\[
\frac1{n!}\sum_{\sigma\in S_n}\chi_{L_n}(\sigma)^2=n,\qquad
\frac1{n!}\sum_{\sigma\in S_n}(\chi_{L_n}(\sigma)-1)^2=n-1.
\]
The report interprets these as top-cohomology invariant multiplicities. In an unnormalized global polynomial sum the associated leading terms are \(nq^n\) and \((n-1)q^n\); averaging over \(q^n\) monic polynomials removes that factor. Thus a raw second moment must not be identified with a centered variance or with its normalized version.

For the short-interval interpretation, intervals are naturally of the form \(I(A;h)=\{f\text{ monic of degree }n:\deg(f-A)\leq h\}\), of size \(q^{h+1}\) for \(0\leq h<n\). Pairing \(L_n\otimes L_n\) with the Hast--Matei top-weight cohomology for pairs of nearby polynomials is said to retain exactly the hooks indexed by
\[
h+2\leq j\leq n-1,
\]
giving coefficient \(n-h-2\) in the Keating--Rudnick short-interval variance. The summary does not state the precise admissible \((n,h)\) range, characteristic conditions, or variance normalization, so no complete asymptotic theorem with a fixed power of \(q\) is supplied for this interpretation.

The candidate synthesis therefore distinguishes three truncations of the same virtual hook object: all hooks give \(n\), all nontrivial hooks give \(n-1\), and the stated short-interval hooks give \(n-h-2\). It does not uniquely determine which expression Ellenberg intended by the incomplete original variance notation.
\subsection{Short English statement}
The original variance expression leaves its exponent and normalization unspecified. The report proposes a geometric hook-representation explanation for three distinct coefficients: a raw moment, a centered global variance, and a short-interval variance.
\subsection{Sources}
\begin{itemize}
\item AIM workshop ``Arithmetic statistics over finite fields and function fields,'' Question 23, attributed to Jordan Ellenberg: \url{https://aimath.org/pastworkshops/arithstatffieldproblems.pdf}. No PDF was accessed here.
\item Local canonical registry \texttt{aim-algebraic-number-theory-notes.json}, \texttt{notes[52]}, tag \texttt{question}; permanent UnsolvedMath ID \(20000401\), source code \texttt{catalogue entry}. Original asks to describe geometrically the right side of a variance relation with \(q^{?}\). Original remarks and literature field are empty.
\item Supplied research attempt 1 invokes Hast--Matei top-weight cohomology and Keating--Rudnick short-interval variance, and proposes the virtual-hook synthesis. No full authoritative bibliographic references or complete asymptotic hypotheses are supplied. The report is machine generated, not peer reviewed, and requires independent verification before citation.
\end{itemize}
\end{document}
